\documentclass[11pt]{article}
\usepackage[margin=0.80in]{geometry}
\usepackage[T1]{fontenc}
\usepackage{lmodern,amsmath,amssymb,amsthm,booktabs,array,xcolor,hyperref}
\hypersetup{colorlinks=true,urlcolor=blue!55!black,linkcolor=blue!45!black}
\newtheorem{theorem}{Theorem}
\newtheorem{lemma}{Lemma}
\newtheorem{proposition}{Proposition}
\newcommand{\CC}{\mathrm C}
\newcommand{\JJ}{\mathrm J}
\newcommand{\QQ}{\mathrm Q}
\newcommand{\KK}{\mathrm K}
\renewcommand{\AA}{\mathrm A}
\newcommand{\word}[1]{\texttt{#1}}
\setlength{\parskip}{5pt}
\setlength{\parindent}{0pt}
\emergencystretch=2em
\title{The minimum trick period in\newline Beggar-My-Neighbour is six}
\author{Christopher D. Long}
\date{18 November 2026}
\begin{document}
\maketitle

\begin{abstract}
For chronological pickup and winner-led play, the least possible primitive
period of a standard 52-card game, measured at completed-trick boundaries and
identifying cards with the same game behaviour, is six. An explicit cycle uses
38 individual card plays and starts at a $30$--$22$ split. A $26$--$26$ deal
enters it after one trick, and the cycle has 181 backward-reachable rank
states, including 19 balanced starts. The upper bound has a short, directly checkable word
certificate. The lower bound reduces all possible periods below six to 1,201
face-card itineraries, followed by exact integer gap feasibility. We also give
an exact classification of the possible deck sizes for period four and explain
why fixed-period deck-size sets are semilinear.
\end{abstract}

\section{Introduction and conventions}

The classical card game Beggar-My-Neighbour (also called Beggar-Your-Neighbor)
was studied as a problem of nontermination for many decades; see the problem
statement of Paulhus~\cite{paulhus} and the historical account of
Casella et al.~\cite{casella}. Casella and collaborators constructed a
nonterminating standard-deck game with a 62-trick cycle, solving Conway's
anti-Hilbert problem. More recent computational and dynamical investigations
include Andorno et al.~\cite{andorno}. The aim here is different: to minimize
the primitive cycle period, under the same deterministic rules and standard
deck inventory. The construction and lower-bound computation described below
show that the minimum is six tricks. We distinguish a cycle's primitive period
from the transient needed to enter it and from a period that tracks the suits
or physical identities of all cards.

This paper is a computationally assisted research manuscript. Its claim of
optimality depends on an exhaustive integer-constraint reduction, accompanied
by reproducible source code and two solver interfaces. The upper-bound witness
is an exact hand-checkable word identity.


Write $\CC$ for any ordinary card and use $\JJ,\QQ,\KK,\AA$ for cards requiring
payments of $1,2,3,4$, respectively. A new face card interrupts the current
payment and starts its own payment. A successful claimant appends the entire
table pile to the bottom of their hand in chronological play order and leads
the next trick. Hands are queues, with the next card on the left.

A state here is $(H_A,H_B,\ell)$ at an empty-table boundary immediately after
collection: both ordered words and the next leader are included. Equal face
ranks are identified, as are all ordinary cards. Let $T$ be the completed-trick
map on states for which another trick can be completed. The primitive period
is the least positive $p$ with $T^p(s)=s$.

\begin{theorem}[Minimum over all splits]
For inventory $\CC^{36}\JJ^4\QQ^4\KK^4\AA^4$,
\[
\boxed{\min_{s\text{ periodic}}\operatorname{per}(s)=6.}
\]
The minimum is attained by a $30$--$22$ state. It is also the minimum possible
eventual period of an infinite game begun with a $26$--$26$ deal.
\end{theorem}

The lower bound below is a computer-assisted exhaustive proof using integer
arithmetic. The existence proof is the explicit six-trick calculation in the
next section. This theorem concerns the minimum over all splits, rather than
a separate function assigning a minimum to each prescribed boundary split.
It also does not assert a minimum for individual card plays or for periods
that distinguish all 52 physical cards.

\section{A six-trick cycle with a word certificate}

Player A leads from the following hands:
\[
\begin{aligned}
H_A={}&\KK\CC\AA\CC^5\QQ\CC^3\QQ\CC\AA\CC^3\QQ
        \KK\CC\AA\CC^5\QQ\CC^2,\\
H_B={}&\CC\JJ\CC\KK\CC\JJ\CC^2\AA\KK\CC^5\JJ\CC^4\JJ\CC.
\end{aligned}
\]
For copying into the verifier, the expanded words are
\begin{verbatim}
A: KCACCCCCQCCCQCACCCQKCACCCCCQCC
B: CJCKCJCCAKCCCCCJCCCCJC
\end{verbatim}
Their lengths are 30 and 22. Together they contain 36 ordinary cards and four
copies of each face rank.

\subsection{The six tricks}

In the table, a pile is written in chronological order. The two play counts
record how many cards each player contributes to that pile.
\begin{center}
\small
\begin{tabular}{@{}c c r r l c c@{}}
\toprule
Trick & Leader & A plays & B plays & Collected pile & Winner & Next split\\
\midrule
1&A&2&2&$P_1=\KK\CC\JJ\CC$&B&$28$--$24$\\
2&B&4&2&$P_2=\CC\AA\KK\CC^3$&B&$24$--$28$\\
3&B&2&2&$P_3=\CC^2\JJ\CC$&B&$22$--$30$\\
4&B&7&8&$P_4=\CC\QQ\CC\AA\CC^3\QQ\KK\CC\AA\CC^4$&A&$30$--$22$\\
5&A&3&2&$P_5=\CC^3\JJ\CC$&B&$27$--$25$\\
6&B&1&3&$P_6=\CC\QQ\CC^2$&A&$30$--$22$\\
\bottomrule
\end{tabular}
\end{center}

For example, the fourth trick has the following sequence of face cards:
player A's queen, B's ace, A's queen, B's king, A's ace. The gaps between
them are $1,3,0,1$, so every interruption is legal; B then pays four ordinary
cards. The complete card-by-card trace is in the accompanying JSON and is
also printed by the verifier.

\subsection{Why both queues return}

Set
\[
U=\KK\CC\AA\CC^5\QQ\CC^2,\qquad V=\CC\JJ\CC.
\]
The starting words have the useful factorizations
\[
\boxed{H_A=UP_4P_6,\qquad H_B=VP_1P_2P_3P_5.}
\]
During the six tricks each player plays exactly 19 cards. No collected card
has yet reached the front of either queue: A starts with 30 cards and B with
22. Their unplayed original suffixes are precisely $U$ and $V$. Their captured
piles are appended in the orders shown in the boxed identities. Consequently
\[
T^6(H_A,H_B,A)=(UP_4P_6,\ VP_1P_2P_3P_5,\ A)
             =(H_A,H_B,A).
\]
The pile lengths are $4,6,4,15,5,4$, totaling 38 card plays. The winner word
$BBBABA$ has least period six. A shorter state period dividing six would force
the winner word to have that shorter period, so the displayed state has
primitive period exactly six.

\section{An equal deal enters the cycle after one trick}

The following two words both have length 26; A leads:
\begin{verbatim}
A: ACCCQKCACCCCCQCCCQCACCCQKC
B: CCCCCJCCCCJCKCJCCAKCCCCCJC
\end{verbatim}
A plays an ace, B pays four ordinary cards, and A collects
$\AA\CC^4$. The resulting hands are
\begin{verbatim}
A: CCCQKCACCCCCQCCCQCACCCQKCACCCC
B: CJCCCCJCKCJCCAKCCCCCJC
\end{verbatim}
with A again leading. This is exactly the boundary following trick 4 in the
six-cycle. Thus the equal deal has a one-trick transient followed by primitive
period six. It is not necessary for a $26$--$26$ boundary itself to occur on
the periodic orbit.

\section{Word equations for a returning FIFO state}

The central combinatorial simplification is independent of the card rules.
Suppose one player begins with a queue $S$, plays a chronological word $D$
during a proposed period, and collects a concatenated word $R$ of captured
piles during that period. If no collected cards have reached the front yet,
then the return equation is $SR=DS$. More generally, the same equation holds
as a concatenation of the chronological queue flow when $D$ is taken to be
the entire word removed from the queue during that interval.

\begin{lemma}[Queue return equation]
Let $S,D,R$ be finite words with $D$ nonempty and $|D|=|R|$. Then
\[
SR=DS\quad\Longleftrightarrow\quad
S=\operatorname{pref}_{|S|}(D^\infty),\qquad
R=\operatorname{rot}_{\,|S|\bmod |D|}(D).
\]
\end{lemma}
\begin{proof}
The equality $SR=DS$ says that $D$ is a prefix of $S R$ and that deleting
this prefix and adding $R$ returns to $S$. Iterating the word equality yields
that $S$ is a prefix of the infinite periodic word $D^\infty$. More
explicitly, write $|S|=q|D|+r$ and $D=XY$ with $|X|=r$. Necessarily
$S=D^qX$ and $R=YX$; conversely, then
$SR=D^qXYX=D^{q+1}X=DS$.
\end{proof}

This lemma imposes only FIFO closure; legality of the intervening face-card
challenge sequence must still be verified. In our six-trick witness, each
player plays 19 cards, and their original queues are longer than 19. The
packet identities in the preceding section instantiate the equation without
replaying collected cards during the cycle.

\section{Backward accessibility from a balanced deal}

For the six-trick orbit, the complete inverse image under the completed-trick
map can be enumerated directly, without a solver. Given a current state and
its most recent winner, the most recent captured pile is a suffix of the
winner's hand. For each possible nonempty suffix and either preceding leader,
we recover the unique prospective ownership sequence from the face-payment
rules; a candidate predecessor is accepted only if one forward trick
reproduces the current state. Iterating these operations from all six
cycle states gives the entire backward component.

\begin{proposition}[Backward component of the displayed orbit]
The six-trick orbit has exactly 181 backward-reachable rank states,
including the six periodic states. Exactly 19 have a 26--26 hand split.
The maximum distance from a backward-reachable state to the cycle is four
completed tricks. Among states with player A leading, the A-hand sizes include
every integer from 20 through 33.
\end{proposition}

The proposition is certified by the standard-library script
\word{backward\_component.py}, which reconstructs all predecessors and
forward-replays each one. Combining it with the global six-trick lower bound
shows that an eventual period of six can be attained from all A-leading
splits $20$--$32$ through $33$--$19$, and the explicit balanced deal above
enters the cycle in one trick. The script also records the complete list of
19 balanced starts for independent inspection. Here the finite counts are
specific to the displayed six-cycle; they are not claims about every
six-trick cycle.

\section{Why the lower-bound search is exhaustive}

\begin{lemma}[Faces and winners]
On a periodic orbit, both players win at least once, and each hand contains
a face card at every completed-trick boundary. Within a trick the owners of
successive face cards alternate. The final face card belongs to the winner.
\end{lemma}
\begin{proof}
If a player never wins, their finite hand loses cards and cannot return.
If a hand ever has no face card, it cannot win another pile and hence cannot
acquire another face card. It therefore cannot lie on a periodic orbit.
After a face is played, the other player must answer; any interrupting face
is consequently owned by the other player. The last face makes the successful
claim when the ensuing payment is completed.
\end{proof}

Suppose trick $i$ has winner $w_i$ and $m_i=2k_i+e_i$ face cards, with
$e_i\in\{0,1\}$. Alternation implies that the loser contributes $k_i$ faces,
and the winner contributes $k_i+e_i$. Hence the winner gains exactly $k_i$
faces. The loser must retain a face; no face in the pile can be used again
before collection. Thus
\[
1\le m_i\le15,\qquad 0\le k_i\le7.
\]
Face balance over a return requires
\[
\sum_{i:w_i=A} k_i=\sum_{i:w_i=B}k_i.
\]

For fixed $p$, the exhaustive enumeration is therefore finite and small:
\begin{enumerate}
\item Enumerate nonconstant binary winner words of length $p$, up to cyclic
rotation and exchange of the players. Choose a representative ending in A,
so the initial leader is A.
\item Enumerate $k_i\in\{0,\ldots,7\}$ and $e_i\in\{0,1\}$ satisfying
the displayed bounds and face balance.
\item Try every initial A-face count $a=1,\ldots,15$. Give the 16 faces
distinct temporary labels and simulate their FIFO movement. Given $w_i$ and
$m_i$, the owner of face $j$ of trick $i$ is forced by alternation. Reject
impossible queues and any boundary with no face in one hand.
\item At a possible return, flatten the two face queues into 16 slots. This
gives a permutation of the initially labelled faces. Every permutation cycle
must be monochromatic in rank. Its cycle sizes must therefore fit into four
bins, each of total weight four.
\item For every surviving face itinerary, solve the exact ordinary-gap
constraints described next, with total ordinary count 36.
\end{enumerate}

This covers every initial card split. It enumerates the number of faces in
each initial hand, but puts no restriction on the number of ordinary cards
assigned to either hand beyond nonnegativity and their combined inventory.

\section{Exact ordinary-gap constraints}

With $a$ faces in A's hand, write its word as
\[
\CC^{g_0}f_1\CC^{g_1}f_2\cdots f_a\CC^{g_a},
\]
and similarly for B. The two hands together have exactly $16+2=18$ gaps.
All initial gaps are nonnegative integers. For the standard deck their sum
is 36. The gap model implements the following three operations exactly.

\paragraph{Before the first face.}
Let $x$ be the current leading gap of the player who will expose the first
face. Let $\delta=0$ if this player is the leader and $\delta=1$ otherwise.
The other player's leading gap must be at least $x+\delta$; these ordinary
cards are consumed. The pile begins with $2x+\delta$ ordinary cards.

\paragraph{Interruptions and final payment.}
After a face of penalty $r$, a prescribed interrupting face can occur after
a leading gap $g$ exactly when $0\le g<r$. The gap is consumed and placed
between the two faces in the pile. If no interruption is prescribed, the
loser's leading gap must be at least $r$; exactly $r$ cards are consumed and
the pile ends in $\CC^r$.

\paragraph{Collection and closure.}
Append the pile to its winner, merging its first ordinary gap with the
winner's existing trailing gap. Carry on with the winner as leader. At the
end, require all 18 gaps to equal their original values. Rank equality is
already imposed by giving each orbit of the face permutation one rank.

These are linear integer constraints, including when the rank attached to
a permutation orbit is a variable in $\{1,2,3,4\}$. For each rank $r$, the
sum of the sizes of its assigned face orbits is constrained to equal four.
The formulation is both necessary and sufficient for a legal return along
the prescribed face itinerary. There is no floating-point tolerance or
sampling in this calculation.

\subsection{The finite certificate for periods below six}

\begin{center}
\begin{tabular}{@{}rrrr@{}}
\toprule
Candidate period & Face itineraries & Rank assignments represented & Feasible at 36 Cs\\
\midrule
2 & 2 & 48 & 0\\
3 & 10 & 288 & 0\\
4 & 163 & 20,328 & 0\\
5 & 1,026 & 134,592 & 0\\
\midrule
Total & 1,201 & 155,256 & 0\\
\bottomrule
\end{tabular}
\end{center}

A one-trick return is impossible because its winner strictly gains cards.
The table excludes every possible return of length two through five. Together
with the explicit six-cycle, this proves the theorem. The table counts
itineraries after symmetry and face-inventory filtering, not all physical
deals and not distinct periodic orbits. The standalone checker regenerates
the table from scratch; no precomputed candidate list is required.

\section{Stronger short-period structure}

The same calculation can leave the number of ordinary cards unrestricted.
Periods two, three and five are impossible for \emph{every} ordinary-card
count when there are four cards of each face rank. Period four has an exact
description.

\subsection{Elementary obstructions for periods two and three}

For period two, only two face itineraries survive the inventory test. Up to
shifting the initial boundary and exchanging the players, the first has
nine A-faces and trick face counts $(8,9)$. If $a_0,b_0$ are initial leading
gaps, B's leading gap returns as $b_0-a_0-r$ for a positive face penalty $r$.
Closure would require $a_0+r=0$, which is impossible. The second itinerary
is the same obstruction with eleven A-faces and counts $(9,8)$.

For period three the winner representative is $BBA$. The ten admissible
itineraries have the following immediate gap contradictions. Initially label
the A-faces $0,\ldots,a-1$ and B-faces $a,\ldots,15$; $r_j>0$ is the penalty
of label $j$. The symbol $\Delta$ denotes final minus initial gap.
\begin{center}
\small
\begin{tabular}{@{}r p{2.35in} l@{}}
\toprule
A-faces & Face counts in the three tricks & Required gap difference\\
\midrule
11 & $(1,8,8),(3,6,8),(5,4,8),(7,2,8)$ & $\Delta a_4=-2b_0-b_5-r_{15}-1<0$\\
9 & $(2,6,9),(4,4,9),(6,2,9)$ & $\Delta b_0=-a_0-r_8<0$\\
13 & $(4,8,13),(6,6,13)$ & $\Delta b_0=-a_0-r_{12}<0$\\
11 & $(7,1,7)$ & $\Delta b_4=b_4+r_{14}>0$\\
\bottomrule
\end{tabular}
\end{center}
Each contradicts the required return of that gap.

\subsection{An exact classification for period four}

\begin{proposition}
With four copies of each face rank and $c$ ordinary cards, a four-trick cycle
exists if and only if
\[
\boxed{c=104+31n,\qquad n\in\mathbb Z_{\ge0}.}
\]
\end{proposition}

The complete face enumeration and unrestricted gap check leave only one
four-trick itinerary, up to the symmetries already used: winners $BABA$, one
face per trick, and initial face words $(rs)^4$ and $(tu)^4$. Here
$(r,s,t,u)$ is any permutation of the four penalties. This finite exclusion
is computer-assisted. The remaining family can be solved algebraically.

Let $q$ be B's initial leading ordinary gap. Define
\[
\begin{aligned}
E&=2q-s+2t+1,& O&=2q+r-2s+2u+1,\\
D&=2q+2r-2s+t+1,& F&=2q+u+1.
\end{aligned}
\]
Return forces the two gap vectors to be
\[
g_A=(E,O,E,O,E,O,E,O,s),\quad
g_B=(q,D,F,D,F,D,F,D,u).
\]
All legal-play inequalities reduce to
\[
q\ge0,\qquad q+r\ge s,\quad q+t\ge s,\quad q+u\ge s.
\]
Since $(r,s,t,u)$ permutes $1,2,3,4$, these are equivalent to $q\ge s-1$.
Adding the gaps gives
\[
c=135+31(q-s)=104+31n,\qquad n=q-s+1\ge0.
\]
Conversely the displayed gap words play as prescribed for every $n\ge0$
and every rank permutation. Their period is four, since period two is
impossible for this face inventory. In particular, the standard count
$c=36$ is excluded by a substantial margin.

\section{The six-trick pattern is rigid}

Suppress ordinary cards in the six-cycle and write the face words as
\[
H_A:\ xyzzyzxyz,\qquad H_B:\ wxwyxww.
\]
The trick face counts are $(2,2,1,5,1,1)$, with winners $BBBABA$. The
physical-face return permutation has four cycles of length four. Thus
$(x,y,z,w)$ must permute $(1,2,3,4)$ if the inventory has four of each rank.

For this particular itinerary, \emph{without specifying an ordinary-card
count}, there is only one legal periodic realization: $(x,y,z,w)=(3,4,2,1)$
and the 36 ordinary cards displayed above. The accompanying
\word{check\_rigidity.py} verifies this by exact rational arithmetic and
ordinary forward play, without an SMT solver.

Here is a finite bound underlying that check. Let $u=b_6$ be B's gap before
its last face. Linear return equations express every gap in $u,x,y,z,w$;
in particular
\[
a_0=\frac{5u-5w+2x-2y-8z+3}{12}.
\]
The second trick requires $a_0<y$, so
\[
5u<5w-2x+14y+8z-3\le85.
\]
Consequently $u\in\{0,\ldots,16\}$. All $24\cdot17$ possibilities can
be checked with rational arithmetic. Only 21 produce nonnegative integral
gap vectors; exactly one has legal play with the prescribed return.

For the successful ranks, all nonnegative integral solutions of the return
\emph{equalities} have the form
\[
\begin{aligned}
g_A={}&(5n,1+5n,5+8n,3+10n,1+5n,\\
      &\quad 3+10n,5n,1+5n,5+8n,2),\\
g_B={}&(1+8n,1+10n,1+3n,2+10n,\\
      &\quad 5n,5+6n,4+12n,1),\qquad n\ge0.
\end{aligned}
\]
Their total ordinary count is $36+115n$. Already the first trick requires
B's jack to interrupt A's king after $b_0-a_0=1+3n$ ordinary cards. Thus
$1+3n<3$, forcing $n=0$. This uniqueness is specific to this itinerary; it
is not a classification of all six-trick cycles.

\section{General structural consequences}

\subsection{Fixed-period deck sizes form a semilinear set}

\begin{theorem}
Fix a finite face-card inventory and a positive integer $p$. The set of
ordinary-card counts for which a completed-trick rank state satisfies
$T^p(s)=s$ is effectively semilinear, and hence eventually periodic. The
same holds for counts admitting a state of primitive period exactly $p$.
More generally, the pairs (ordinary-card count, initial A-hand size) form an
effectively semilinear subset of $\mathbb N^2$.
\end{theorem}
\begin{proof}
Let $f$ be the fixed number of face cards. Both players retain a face at each
periodic boundary, so a trick contains at most $f-1$ faces. For fixed $p$
there are only finitely many possible face itineraries and rank assignments.
Each is described by linear integer conditions on its $f+2$ ordinary gaps.
Their finite union is therefore definable in Presburger arithmetic. Adding
the equations for total ordinary count and initial A-hand size, then
projecting out the gaps, preserves Presburger definability. Presburger sets
are effectively semilinear \cite[Theorem 1.3]{ginsburg-spanier}.

For primitive period $p$, exclude every proper-divisor return on the
\emph{same initial state variables before projection}. Along a fixed
itinerary, such a return is either impossible because of the leader or
face word, or is expressed by linear gap equalities. Negating these
equalities remains Presburger-definable. One must not subtract sets of
ordinary counts after projection: one count can support cycles of more
than one primitive period.
\end{proof}

The progression $104+31n$ for period four is a concrete example. The result
also explains why an exact search based on gaps can classify families with
arbitrarily many ordinary cards, rather than requiring a finite deck-size
cutoff.

\subsection{Lifting a rank cycle to individually labelled cards}

If a rank state returns after $p$ tricks, give each card a distinct label and
follow those same $p$ tricks. The resulting card-slot permutation preserves
the game-relevant types. Repeating the rank itinerary applies the same
permutation again. Hence the labelled period equals $p$ times the order of
this permutation when $p$ is the primitive rank period.

For the displayed six-cycle the face-label cycles have lengths $4,4,4,4$;
the ordinary-label cycles have lengths $3,3,3,3,4,9,11$. Their least common
multiple is $396$, so the labelled-card period of this example is
\[
6\cdot396=2376\text{ tricks},
\]
or $396\cdot38=15{,}048$ individual plays. This does not make 2376 a claimed
minimum among labelled-card cycles.

\section{Reproduction and scope of verification}

The portable script \word{prove\_minimum.py} requires Python 3 and
\word{z3-solver}. It regenerates all face itineraries for periods two
through five, checks their exact integer constraints, verifies the explicit
six-cycle with two implementations of play, and checks the balanced entry.
It raises an error on any solver result other than \word{unsat} in the
lower-bound checks. No timeout or unresolved case is counted as an exclusion.
\begin{verbatim}
python -m pip install z3-solver
python prove_minimum.py --json minimum_verification.json
python prove_minimum.py --structural
python check_rigidity.py --json rigidity_verification.json
\end{verbatim}
The optional \word{--smt-dir DIRECTORY} argument exports every ordinary-gap
problem in SMT-LIB format for inspection or another solver. The optional
structural checks remove the ordinary-count restriction for periods two,
three and five and verify the period-four classification.

A second backend, \word{native\_z3\_check.py}, reuses the same exact gap
models but serializes them directly as SMT-LIB and calls the native Z3 C
library through Python's \word{ctypes}. It does not require the Python
\word{z3-solver} package. A fresh run with native Z3 4.13.3 returned
\word{unsat} for all 1,201 candidate face itineraries for periods two
through five. This verifies the feasibility computations using a distinct
solver interface, but it does not replace the need to audit the underlying
model reduction.

The backward component can be independently reproduced without third-party
Python packages:
\begin{verbatim}
python backward_component.py --json verification/backward_component.json
python native_z3_check.py --json verification/native_z3.json
\end{verbatim}

The minimum proof was crosschecked using independently written gap
bookkeeping, an independent winner-necklace enumerator, and two direct
game simulators. The fixed-period exclusions are still a computational
certificate, not a proof-assistant formalization. The explicit upper-bound
and balanced-entry calculations can be checked directly from the words.

No claim is made here that the six-cycle is the unique cycle of minimum
period, that 38 is the least possible number of individual plays, or that
every prescribed split occurs on a six-cycle. What is established is the
minimum completed-trick rank period over all splits, and its attainment as
an eventual period from an equal deal.

\begin{thebibliography}{9}
\bibitem{paulhus}
M. M. Paulhus,
\emph{Beggar My Neighbour},
American Mathematical Monthly \textbf{106} (1999), 162--165.
\href{https://doi.org/10.2307/2589054}{doi:10.2307/2589054}.
\bibitem{casella}
B. Casella, P. M. Anderson, M. Kleber, R. P. Mann, R. Nessler, W. Rucklidge, S. G. Williams, and N. Wu,
\emph{A non-terminating game of Beggar-My-Neighbor},
American Mathematical Monthly \textbf{132}(10) (2025), 978--994.
\href{https://doi.org/10.1080/00029890.2025.2558484}{doi:10.1080/00029890.2025.2558484}.
\bibitem{andorno}
N. Andorno, G. Cernoia, S. Duiz, and A. Michelangeli,
\emph{The dynamical landscape of Beggar-My-Neighbour: Ultra-long matches, loops, and infinite matches},
arXiv:2602.23406 (2026).
\href{https://arxiv.org/abs/2602.23406}{arxiv.org/abs/2602.23406}.
\bibitem{ginsburg-spanier}
S.~Ginsburg and E.~H.~Spanier,
\emph{Semigroups, Presburger formulas, and languages},
Pacific Journal of Mathematics \textbf{16}(2) (1966), 285--296.
\href{https://doi.org/10.2140/pjm.1966.16.285}{doi:10.2140/pjm.1966.16.285}.
\end{thebibliography}
\end{document}
